\documentclass[12pt,a4paper]{article}
\usepackage[top=2.5cm,bottom=2.5cm,left=3cm,right=2.5cm]{geometry}
\usepackage[T1]{fontenc}
\usepackage{newtxtext,newtxmath}
\usepackage[table]{xcolor}
\usepackage{graphicx,booktabs,amsmath,enumitem,float,caption,setspace,fancyhdr,titlesec}
\usepackage[colorlinks=true,linkcolor=black,urlcolor=blue!50!black,citecolor=black]{hyperref}
\onehalfspacing
\definecolor{bench}{RGB}{226,239,218}
\definecolor{chosen}{RGB}{221,235,247}
\setlist{itemsep=1pt,topsep=3pt}
\captionsetup{font=small,labelfont=bf}
\titleformat{\section}{\large\bfseries}{\thesection.}{0.6em}{\MakeUppercase}
\titleformat{\subsection}{\normalsize\bfseries}{\thesubsection}{0.6em}{}
\pagestyle{fancy}
\fancyhf{}
\fancyhead[L]{\small UCS760 Course Project}
\fancyhead[R]{\small DQN Traffic Signal Control, Pune}
\fancyfoot[C]{\small\thepage}
\renewcommand{\headrulewidth}{0.4pt}
\newcommand{\trainsteps}{150,000}
\newcommand{\trainenvs}{4}
\newcommand{\trainseeds}{2}
\newcommand{\camseedsword}{two}
\newcommand{\routeseeds}{2}
\newcommand{\routeseedsword}{two}
\newcommand{\trainminutes}{24--39}
\newcommand{\fig}[2]{\IfFileExists{#2}{\includegraphics[width=#1]{#2}}{\fbox{\parbox{#1}{\centering\vspace{1.5cm}\textit{Figure pending: results not generated yet}\vspace{1.5cm}}}}}

\newcommand{\student}{Mridul Negi}
\newcommand{\rollno}{102303118}
\newcommand{\teacher}{Himika Sharma}
\newcommand{\repourl}{\url{https://github.com/MridulNegi2005/dqn-traffic-signal-pune}}

\begin{document}
\begin{titlepage}
\centering
\vspace*{1.5cm}
{\large\bfseries THAPAR INSTITUTE OF ENGINEERING AND TECHNOLOGY, PATIALA\par}
\vspace{0.3cm}
{\normalsize Department of Computer Science and Engineering\par}
\vspace{1.6cm}
{\large Project Report\par}
\vspace{0.5cm}
{\LARGE\bfseries Deep Q-Network Traffic Signal Control on Real Heterogeneous Traffic from Pune\par}
\vspace{1.2cm}
{\large UCS760: Edge AI and Robotics --- Reinforcement Learning\par}
\vfill
\begin{minipage}[t]{0.45\textwidth}
\raggedright\textbf{Submitted by}\\[2pt]
\student\\
Roll No.\ \rollno
\end{minipage}\hfill
\begin{minipage}[t]{0.45\textwidth}
\raggedleft\textbf{Submitted to}\\[2pt]
\teacher
\end{minipage}
\vfill
\includegraphics[width=0.32\textwidth]{tiet_logo.jpg}\par
\vspace{0.4cm}
{\large November 2026\par}
\vspace*{1cm}
\end{titlepage}

\section*{Abstract}
\addcontentsline{toc}{section}{Abstract}

Most reinforcement learning (RL) studies of traffic signal control use synthetic, car-only demand
and lane-disciplined traffic. Indian junctions carry mixed traffic in which motorbikes are close to
half of all vehicles, and many junctions serve one road at a time. This project trains a Deep
Q-Network (DQN) agent that decides which road of a four-way junction gets the green light and for how
long. The demand comes from real CCTV vehicle counts at Jehangir Chowk, Pune (11--18 January 2023).
The SUMO simulation models four vehicle classes, left-hand traffic, dedicated turn lanes and lane
sharing by motorbikes; the reward weights each vehicle by its Passenger Car Unit (PCU). The agent is
trained on five days with randomised demand and tested on three unseen days, on 1.3 times the real
demand, on a junction it never saw (RTO Chowk) and with free left turns. It is compared with a
fixed-time plan, a loop-sensor actuated controller and a hand-written camera rule. All adaptive
controllers cut the mean waiting time by about 80\% against the fixed plan. The camera rule,
which needs only a CCTV camera, waits 25--43\% less than the actuated controller in normal traffic;
the camera-view DQN, warm-started from this rule, comes close to the actuated controller.


\clearpage
\section{Introduction}
A traffic signal decides which approach of a junction gets the green light and for how long.
Many Indian junctions run a fixed plan that gives each road the same green time in a fixed order.
A fixed plan does not react to the queues that are present. Reinforcement learning lets a controller
learn a policy by interaction with a simulator: it observes the queues, selects a road to serve, and
receives a reward that measures the delay of the vehicles.

Most Indian junctions have CCTV cameras but no induction loops in the road. This project
therefore asks what a signal controller can do with the view of one camera at the stop line.
The contributions are:
\begin{enumerate}
  \item A pipeline that converts the per-frame running totals of the Pune CCTV dataset into
        5-minute, per-class arrival rates for SUMO.
  \item A Gymnasium environment for a left-hand, mixed-traffic junction with one-road-at-a-time
        phasing, dedicated left, straight and right lanes, an optional free left turn, and a
        simulated camera that sees only the last 75\,m of each road.
  \item A combined reward that counts queues and spill-back in PCU and charges for each change of
        road, and a camera-only rule that serves as a baseline and as a teacher for DQN.
  \item A comparison of ten training strategies and a final test on unseen days, heavier demand, an
        unseen junction and free left turns, with five training seeds, against fixed-time and
        loop-sensor actuated control.
\end{enumerate}

\section{Related Work}
Genders and Razavi showed that a deep Q-network can control a single signalised
intersection in SUMO and reduce delay compared with a fixed plan. IntelliLight
trained a DQN on real traffic volumes from cameras in Jinan, China, and tested it in a simulator.
Wei et al. survey later work on multi-intersection control and state and reward
design. Most of these studies use lane-disciplined, car-only traffic and right-hand driving.
This project uses real Indian counts with four vehicle classes, left-hand driving, motorbikes that
share lanes, a PCU-weighted reward, and the one-road-at-a-time phasing that is common in India.

\section{Background}
\subsection{Markov decision process}
The control problem is a Markov decision process (MDP) $(\mathcal{S},\mathcal{A},P,R,\gamma)$.
At each decision step $t$ the agent observes a state $s_t$, selects an action $a_t$,
receives a reward $r_{t+1}$ and moves to the state $s_{t+1}$. The agent maximises the expected
discounted return $G_t=\sum_{k\ge 0}\gamma^k r_{t+k+1}$.

\subsection{Q-learning and DQN}
The action-value function $Q^\pi(s,a)$ is the expected return after action $a$ in state $s$.
Q-learning is a temporal-difference (TD) method that moves its estimate towards the Bellman
optimality target:
\begin{equation}
Q(s_t,a_t) \leftarrow Q(s_t,a_t) + \alpha\Big[r_{t+1} + \gamma \max_{a'} Q(s_{t+1},a') - Q(s_t,a_t)\Big].
\end{equation}
A table of $Q$ values is not possible for a continuous state such as queue lengths.
DQN replaces the table with a neural network $Q_\theta$ and adds two parts
for stable training:
\begin{itemize}
  \item an \emph{experience replay buffer} that stores transitions and samples random mini-batches;
  \item a \emph{target network} $Q_{\theta^-}$, a delayed copy of $Q_\theta$ that gives the target
        $y=r+\gamma\max_{a'}Q_{\theta^-}(s',a')$.
\end{itemize}
The loss is $\big(y-Q_\theta(s,a)\big)^2$. During training, an $\varepsilon$-greedy policy selects
a random action with probability $\varepsilon$, so the agent tries choices it would not otherwise
make and learns from the result. During testing the agent always takes the action with the highest
$Q$ value.

\subsection{Passenger car unit}
A PCU converts a vehicle of any class into an equivalent number of passenger cars.
IRC:106-1990 gives 1.0 for a car, 0.5 for a motorbike and 2.2 for a bus or a truck.

\section{Dataset}
The \emph{Heterogeneous Traffic Count Dataset, Pune (Jan 2023)} contains vehicle
counts from CCTV at three signalised junctions in Pune. The Traffic Police Commissioner, Pune
supplied the video. A YOLO detector and a multi-object tracker produced counts for four classes
(car, motorbike, bus, truck) and four directions of movement in the image.

\subsection{Preprocessing}
Each CSV row is one video frame. The values are running totals since the start of the
recording session; a file with the suffix \texttt{\_one} continues a session and restarts at zero.
The script \texttt{src/data/preprocess.py} takes the last running total in each 5-minute bin and the
difference between consecutive bins. The output has 9{,}528 rows: one row per junction, camera,
day, 5-minute bin and direction. All cameras cover 09:00--12:00 on each of the eight days.

\subsection{Demand for the simulator}
Jehangir Chowk and RTO Chowk each have three cameras. Each camera is mapped to the north, east or
south road. Each camera sees traffic in both directions, so the inbound rate is half of the
camera count. The west road has no camera; it is a minor road with 40\% of the mean of the other
roads. The dataset has no turning counts; the split is 70\% straight, 15\% left and 15\% right.
For each 5-minute bin and class, departures follow a Poisson process with the observed rate.
The mean inbound demand at Jehangir Chowk is about 3{,}650 vehicles per hour, and almost half of
the vehicles are motorbikes (Figure~\ref{fig:demand}).

\begin{figure}[H]
\centering
\fig{\linewidth}{../results/demand_profile.png}
\caption{Left: inbound demand per road in 15-minute steps, mean of the eight days. Right: share
of each vehicle class.}
\label{fig:demand}
\end{figure}

\section{Simulation Environment}
\subsection{Road network and vehicles}
Figure~\ref{fig:junction} shows the simulated junction.

\begin{figure}[H]
\centering
\fig{0.62\linewidth}{../results/junction.png}
\caption{The simulated junction (top view, left-hand traffic). Arrows show the incoming traffic;
the orange zones are the parts of each road that the camera sees.}
\label{fig:junction}
\end{figure}

The simulator is SUMO. The network is a four-way junction with 250\,m roads, built
with \texttt{netconvert --lefthand}.
Each road has three incoming lanes: a left-turn lane on the kerb side, a straight lane and a
straight-and-right lane. Motorbikes use the SUMO sublane model with a lateral resolution of
0.8\,m, so they can move between cars. Emissions use HBEFA classes: a Euro~4 petrol car, a
4-stroke motorbike up to 250\,cc, a city bus and a heavy diesel truck.

\subsection{Signal phasing}
The signal serves one road at a time (split phasing): all movements of the active road are green
and every other road is red. This is a common plan at Pune junctions, because a protected
right-turn phase is rare and opposing right turns conflict. A change of road always passes through
3\,s of yellow and 2\,s of all-red.

In the \emph{free left} variant, the left-turn lane works as a slip lane: left turns may go on red
but must give way to the active movements.

\clearpage
\section{Methodology}
\subsection{What the controller can see}
A real junction in an Indian city has CCTV cameras, not induction loops in every lane. A camera
with a detector and a tracker (as used to build the Pune dataset) can count vehicles by class,
tell if they move, and measure how long each stopped vehicle has waited, but only within its field
of view. The agent therefore sees only the last 75\,m of each incoming lane.

\subsection{MDP formulation}
\begin{description}[leftmargin=1.6cm,style=nextline]
  \item[State] For each of the 12 incoming lanes: the PCU of halted vehicles and the PCU of all
        vehicles inside the camera view, each divided by the PCU that fits in the view; for each
        road, the longest waiting time of a vehicle in view (divided by 300\,s); a one-hot code of
        the road that has green; a flag that shows if the minimum green has passed; and the
        elapsed green time. The state has 34 values in $[0,1]$.
  \item[Action] The road that gets green next (4 actions). The agent decides every $\Delta t$
        seconds. If it keeps the same road, the green continues. If it selects another road and the
        minimum green (10\,s) has passed, the signal shows 3\,s of yellow and 2\,s of all-red and
        then the new green. A green that reaches the maximum (60\,s) passes to another road.
        In the \emph{keep/switch} variant the agent has 2 actions: keep the green, or switch to the
        road that the camera rule (below) ranks highest.
  \item[Reward] For each second $s$ of a decision step, the controller pays a cost; vehicles that
        cross the stop line earn a bonus; a change of road costs a fixed penalty:
        \begin{equation}
        r_{t+1} = -\sum_{s}\Big(w_q\,Q_s + w_b\,B_s + w_m\,M_s\Big) + w_x\,X_t - w_c\,\,\mathbf{1}[\text{switch}],
        \end{equation}
        where $Q_s$ is the PCU of stopped vehicles on the incoming lanes, $B_s$ is the PCU of
        vehicles that cannot enter the road because the queue has spilled back, $M_s$ is the
        longest accumulated waiting time of any vehicle, and $X_t$ is the PCU that crossed the stop
        line during the step. The weights are $w_q=w_b=0.01$, $w_m=0.002$, $w_x=0.2$ and
        $w_c=0.5$. The terms target queues, spill-back, starvation of a quiet road, throughput and
        the 5\,s lost at each change of road. The reward uses simulator data that a camera cannot
        give (for example $B_s$); this is allowed because the reward is used only in training,
        while the deployed controller uses only the camera state.
\end{description}

\subsection{Camera rule: a baseline and a teacher}
The camera rule is a hand-written controller that uses only camera data. It keeps the green while
vehicles are inside the last 30\,m of the green road, between the minimum and maximum green. When
the green road is clear, it serves the road with the highest score
$\sum_{\text{lanes}} \mathrm{PCU}_{\text{halted}} + 0.05\times(\text{longest wait in s})$,
so it skips empty roads. It is a baseline and also the \emph{teacher} for a warm start: before RL
training, 20{,}000 decisions of the camera rule fill the replay buffer, and the exploration rate
starts at 0.2 instead of 1.0. The agent then learns to improve on the rule.

\subsection{Strategies}
Ten training runs test one change each (Table~\ref{tab:strategies}). All use the camera state and
the combined reward unless the table says otherwise.

\begin{table}[H]
\centering\small
\begin{tabular}{lll}
\toprule
Run & Change & Question \\
\midrule
R1 & DQN, $\Delta t=5$\,s & Effect of the new reward and state alone \\
R2 & DQN, $\Delta t=2$\,s & Effect of faster decisions \\
R3--R5 & R2 + warm start from the camera rule (3 seeds) & Main candidate and its spread \\
R6 & R3 + keep/switch action & Is a 2-action design easier to learn? \\
R7 & R3 with QR-DQN & Effect of a distributional DQN \\
R8 & R3 trained for 300{,}000 steps & Effect of longer training \\
R9 & R3 with the old reward (drop in waiting time) & Contribution of the combined reward \\
R10 & R3 with a full-lane view instead of the camera & Cost of camera-only sensing \\
\bottomrule
\end{tabular}
\caption{Training runs. Each run trains for 150{,}000 steps unless noted.}
\label{tab:strategies}
\end{table}

\subsection{Training and testing}
\begin{itemize}
  \item \textbf{Training:} days 11--15 January at Jehangir Chowk. Each episode is one random hour
        of the morning peak with the demand scaled by a random factor in $[0.7,1.4]$.
        DQN from Stable-Baselines3, MLP with two hidden layers of 256 units,
        learning rate $5\times10^{-4}$, replay buffer $2\times10^5$, batch 128, $\gamma=0.99$,
        target update every 2{,}000 steps. Each run uses 4 parallel SUMO instances on a Kaggle
        notebook.
  \item \textbf{Test scenarios:} days 16--18 January, full 3-hour peak, three route seeds per day:
        (1) unseen days at Jehangir Chowk; (2) the same days at 1.3 times the real demand;
        (3) an unseen junction, RTO Chowk; (4--5) unseen days at 1.0 and 1.3 times demand with
        free left turns. Each route seed fixes the vehicle departures, so all controllers see the
        same vehicles.
  \item \textbf{Baselines:} a fixed-time plan (30\,s green per road, 140\,s cycle); the SUMO
        actuated controller on the same four phases, which uses an induction loop on every lane
        and extends a green between 10 and 60\,s while vehicles arrive; and the camera rule.
  \item \textbf{Metrics:} from SUMO trip information. Waiting time and delay include the time a
        vehicle waits to enter the network. Throughput is the number of vehicles that complete
        their trip per hour. Fuel and CO$_2$ are per vehicle.
\end{itemize}



\clearpage
\section{Results}
\subsection{First attempt}
The first DQN used the drop in waiting time as its reward, a 5\,s decision step and a full-lane
view. On the unseen days it reduced the mean waiting time from 485\,s (fixed-time) to 294\,s, but
the actuated controller reached 98\,s. Its reward did not see vehicles that could not enter the
road, it did not pay for a change of road, and it could leave a quiet road unserved for a long
time. These findings led to the combined reward and the strategies in Table~\ref{tab:strategies}.

\subsection{Strategy comparison}
Table~\ref{tab:ablation} and Figure~\ref{fig:ablation} compare the ten training runs with the
three baselines on the unseen days at 1.0 and 1.3 times the real demand (one route seed per day).

\begin{table}[H]
\centering\small
\begin{tabular}{lrrrrr}
\toprule
 & \multicolumn{2}{c}{Mean wait (s)} & \multicolumn{2}{c}{Throughput (veh/h)} & CO$_2$ \\
\cmidrule(lr){2-3}\cmidrule(lr){4-5}
Controller & 1.0$\times$ & 1.3$\times$ & 1.0$\times$ & 1.3$\times$ & (g/veh) \\
\midrule
\multicolumn{6}{l}{\textit{Baselines}} \\
Fixed-time & 495 & 952 & 3033 & 3459 & 481 \\
Actuated (loop sensors) & 86 & 491 & 3393 & 3968 & 300 \\
\rowcolor{bench}\textbf{Camera rule} & \textbf{66} & \textbf{453} & \textbf{3401} & \textbf{4051} & \textbf{298} \\
\midrule
\multicolumn{6}{l}{\textit{DQN training runs}} \\
R1 DQN, 5 s & 337 & 499 & 2502 & 2704 & 543 \\
R2 DQN, 2 s & 308 & 945 & 2988 & 2916 & 474 \\
\rowcolor{chosen}\textbf{R3--R5 warm start, camera view (3 seeds)} & \textbf{109} & \textbf{622} & \textbf{3371} & \textbf{3816} & \textbf{320} \\
R6 keep/switch action & 217 & 526 & 3283 & 3914 & 358 \\
R7 QR-DQN & 379 & 650 & 3104 & 3788 & 370 \\
R8 300k steps & 431 & 996 & 3097 & 3477 & 429 \\
R9 old reward & 150 & 549 & 3352 & 3889 & 358 \\
R10 full-lane view & 76 & 467 & 3408 & 4010 & 326 \\
\bottomrule
\end{tabular}
\caption{Strategy comparison on unseen days (16--18 January). Green: the camera rule, used as the
benchmark and teacher. Blue: the strategy chosen for the final test (warm start, camera view);
its row is the mean of three seeds.}
\label{tab:ablation}
\end{table}

\begin{figure}[H]
\centering
\fig{0.85\linewidth}{../results/ablation.png}
\caption{Mean waiting time of each strategy and baseline on unseen days.}
\label{fig:ablation}
\end{figure}

\paragraph{What the strategy comparison shows.}
\begin{itemize}
  \item \textbf{Warm start is the largest single gain.} Without it, DQN with the new reward (R2)
        waits 308\,s at real demand and 945\,s at 1.3$\times$ demand, close to the fixed plan.
        With the camera-rule warm start (R3--R5) the mean is 109\,s and 622\,s.
  \item \textbf{A longer view helps the learned controller.} R10 (warm start, full 236\,m view)
        waits 76\,s against 86\,s for actuated control (11\% less) and 467\,s against 491\,s at
        1.3$\times$ demand (5\% less). One camera at the stop line cannot give this view; it needs
        a second camera up the road or a higher mounting, so R10 is an upper bound and not part of
        the camera-only comparison.
  \item \textbf{The hand-written camera rule is the strongest controller in this test}
        (66\,s and 453\,s). It already beats actuated control, because it skips empty roads and
        ranks roads by PCU queue and waiting time. DQN starts from this rule but does not keep all
        of its quality; with a camera-only view the three seeds differ widely (96\,s, 87\,s and
        146\,s at real demand).
  \item \textbf{Longer training is worse, not better.} R8 (300{,}000 steps) waits 431\,s against
        109\,s for the same setup at 150{,}000 steps. The agent drifts away from the teacher's
        policy as training continues, which points to a need for early stopping or a stronger
        anchor to the teacher.
  \item \textbf{The other changes do not help in this setting}: the keep/switch action (R6),
        QR-DQN (R7) and the 5\,s step (R1). R1 has a low waiting time at 1.3$\times$ demand only
        because 37\% of its vehicles never entered the road and so have no waiting time recorded;
        its throughput is the lowest of all controllers.
  \item \textbf{Reward:} the combined reward gives a lower waiting time than the old reward at real
        demand (109\,s against 150\,s, R3--R5 against R9) but a higher one at 1.3$\times$ demand
        (622\,s against 549\,s), mostly because of one weak seed.
\end{itemize}
The full test therefore uses the camera rule and the camera-view DQN (the R3--R5 setup), which
both work with one camera at the stop line. Its result in the full test is the mean of several training seeds.


\subsection{Full test}
The full test compares the camera-view DQN (the R3--R5 setup, trained with \camseedsword{} seeds) with the
three baselines on all five scenarios, with \routeseedsword{} route seeds per day (Table~\ref{tab:summary},
Figure~\ref{fig:comparison}). Every controller in this test uses only data that one camera or one
set of loops at the junction can give.

\begin{table}[H]
\centering\scriptsize
\setlength{\tabcolsep}{3.5pt}
\begin{tabular}{llrrrrrr}
\toprule
Scenario & Controller & Wait (s) & Delay (s) & Trip (s) & Thr.\ (veh/h) & CO$_2$ (g) & Fuel (g) \\
\midrule
Unseen days & Fixed-time & 484.8 & 509.6 & 184.2 & 3024 & 477 & 151 \\
 & Actuated (loops) & 98.0 & 111.1 & 106.9 & 3368 & 298 & 94 \\
\rowcolor{bench} & Camera rule & 63.4 & 76.4 & 106.0 & 3392 & 294 & 93 \\
\rowcolor{chosen} & \textbf{DQN, camera view} & 103.7 $\pm$ 27.5 & 118.1 & 112.5 & 3361 & 310 & 98 \\
\midrule
1.3$\times$ demand & Fixed-time & 970.0 & 1001.6 & 227.6 & 3457 & 568 & 180 \\
 & Actuated (loops) & 504.4 & 522.3 & 163.2 & 3963 & 426 & 135 \\
\rowcolor{bench} & Camera rule & 507.3 & 526.5 & 175.9 & 4026 & 450 & 143 \\
\rowcolor{chosen} & \textbf{DQN, camera view} & 538.8 $\pm$ 71.1 & 559.1 & 177.7 & 3968 & 454 & 144 \\
\midrule
Unseen junction (RTO) & Fixed-time & 535.4 & 556.8 & 161.8 & 2710 & 427 & 135 \\
 & Actuated (loops) & 53.4 & 65.6 & 92.6 & 3100 & 277 & 88 \\
\rowcolor{bench} & Camera rule & 40.2 & 51.5 & 87.3 & 3101 & 262 & 83 \\
\rowcolor{chosen} & \textbf{DQN, camera view} & 48.9 $\pm$ 6.1 & 61.9 & 95.7 & 3100 & 284 & 90 \\
\midrule
Free left & Fixed-time & 503.4 & 528.3 & 181.1 & 3018 & 470 & 149 \\
 & Actuated (loops) & 98.7 & 111.7 & 106.7 & 3359 & 298 & 94 \\
\rowcolor{bench} & Camera rule & 55.8 & 68.2 & 98.8 & 3403 & 278 & 88 \\
\rowcolor{chosen} & \textbf{DQN, camera view} & 239.0 $\pm$ 36.6 & 253.5 & 120.8 & 3176 & 326 & 103 \\
\midrule
Free left, 1.3$\times$ & Fixed-time & 977.1 & 1007.9 & 217.9 & 3445 & 547 & 173 \\
 & Actuated (loops) & 512.9 & 530.5 & 157.3 & 3936 & 413 & 131 \\
\rowcolor{bench} & Camera rule & 431.8 & 450.0 & 159.2 & 4035 & 415 & 132 \\
\rowcolor{chosen} & \textbf{DQN, camera view} & 737.2 $\pm$ 36.8 & 757.1 & 181.6 & 3391 & 458 & 145 \\
\bottomrule
\end{tabular}
\caption{Full test results. Mean over 3 days $\times$ \routeseeds{} route seeds; for DQN also over \trainseeds{} training
seeds ($\pm$ is the standard deviation between training seeds). Thr.\ is throughput; CO$_2$ and
fuel are per vehicle.}
\label{tab:summary}
\end{table}

\begin{figure}[H]
\centering
\fig{\linewidth}{../results/comparison.png}
\caption{Waiting time, trip time, throughput and CO$_2$ for each controller and scenario.
Error bars show the spread between the \camseedsword{} training seeds of the DQN.}
\label{fig:comparison}
\end{figure}

\begin{figure}[H]
\centering
\fig{0.8\linewidth}{../results/learning_curve.png}
\caption{Episode return of the camera-view DQN during training. Thin lines: single training
seeds; thick line: their mean.}
\label{fig:learning}
\end{figure}

\begin{table}[H]
\centering\small
\begin{tabular}{lrrrr}
\toprule
Controller & Car & Motorbike & Bus & Truck \\
\midrule
Fixed-time & 525.2 & 485.2 & 402.2 & 526.4 \\
Actuated (loops) & 111.5 & 103.2 & 73.6 & 103.6 \\
Camera rule & 67.7 & 67.9 & 65.2 & 63.3 \\
\textbf{DQN, camera view} & 116.2 & 110.3 & 81.5 & 106.9 \\
\bottomrule
\end{tabular}
\caption{Mean waiting time (s) by vehicle class on the unseen days.}
\label{tab:class}
\end{table}

\subsection{Discussion}
\begin{itemize}
  \item \textbf{All adaptive controllers are far better than the fixed plan.} On the unseen days
        the fixed plan gives a mean wait of 485\,s; the actuated controller, the camera rule and
        the camera-view DQN give 98\,s, 63\,s and 104\,s. They also move about 11\% more vehicles
        per hour and use about 35\% less fuel per vehicle.
  \item \textbf{The camera rule is the best controller.} With only a 75\,m camera view it waits
        35\% less than the loop-sensor actuated controller on the unseen days (63\,s against
        98\,s), 25\% less on the unseen junction (40\,s against 53\,s) and 43\% less with free
        left turns. At 1.3$\times$ demand it equals the actuated controller (507\,s against 504\,s)
        and has a higher throughput. Its advantage comes from two choices: it skips roads with no
        queue, and it ranks roads by the PCU of the queue and the longest wait.
  \item \textbf{The camera-view DQN is close to the actuated controller.} On the unseen junction it
        waits 49\,s against 53\,s, and on the unseen days 104\,s against 98\,s. The result changes
        between training seeds (standard deviation 28\,s on the unseen days), so the mean of
        several seeds is reported.
  \item \textbf{The DQN does not handle a change of road layout.} With free left turns, which it
        never saw in training, its wait rises to 239\,s, while the two rules do not need training
        and stay good. A learned controller must train on the layout where it will run.
  \item \textbf{Buses and trucks are not left behind.} For every adaptive controller, buses wait
        less than cars and trucks wait about as long as cars (Table~\ref{tab:class}).
\end{itemize}



\section{Limitations and Future Work}
The dataset gives no turning movements and has no camera for the west road, so the turn split
(70/15/15) and the west-road demand are assumptions. The camera counts come from a YOLO detector
and cover only 09:00--12:00 in January. Split phasing serves fewer vehicles than plans in which
opposite roads move together; all controllers use the same phases, so the comparison stays fair.
The simulated camera is perfect: it detects every vehicle in its 75\,m view, while a real
detector misses vehicles under occlusion, at night or in rain, so real results would be somewhat
worse for the camera rule and the camera-view DQN. The DQN results differ between training seeds;
the report gives the mean and the spread of five seeds instead of the best seed.

Future work:
\begin{itemize}
  \item control several connected junctions with multi-agent RL, so that one signal knows what
        its neighbour will send;
  \item add pedestrian phases and protected right-turn arrows;
  \item collect turning counts from the CCTV video to remove the turn-split assumption;
  \item add detection errors to the simulated camera and train the agent to be robust to them;
  \item keep the learned policy close to the teacher (for example with early stopping or a
        penalty for leaving the teacher's action), because longer training made the agent worse;
  \item deploy the trained network on an edge device such as an NVIDIA Jetson next to the camera,
        so that detection and control run on the same unit.
\end{itemize}

\section{Conclusion}
This project built a traffic-signal controller for a Pune junction that uses only the view of a
CCTV camera. Real vehicle counts from Jehangir Chowk drive a SUMO model with mixed traffic,
left-hand driving and one-road-at-a-time signals.

The main results are:
\begin{itemize}
  \item A DQN agent trained from scratch learns to beat the fixed plan but stays far behind the
        actuated controller. A warm start from a simple camera rule, a combined reward and a
        decision every 2\,s close most of that gap.
  \item The camera rule itself is the strongest controller. It waits 25--43\% less than the
        loop-sensor actuated controller in normal traffic, equals it at 1.3 times the demand, and
        needs only a camera that most junctions already have.
  \item The camera-view DQN comes close to the actuated controller and works on a junction it never
        saw, but its result depends on the training seed and it fails when the road layout
        changes.
\end{itemize}

For a single junction, a well-designed camera rule is the practical choice today. Reinforcement
learning is most useful as a way to improve on such a rule, and its value grows when several
junctions must work together, which is the next step for this work.


\clearpage
\appendix
\section{Reproducibility}
All code, the SUMO network and the processed 5-minute counts are in the repository \repourl.
The raw dataset (\url{https://data.mendeley.com/datasets/xnf2k6n288}, DOI 10.17632/xnf2k6n288.1,
CC BY 4.0) is downloaded from Mendeley Data by a script that checks the SHA-256 hash of each file.
\begin{itemize}
  \item \textbf{Software:} Python 3.11+, SUMO 1.28.0 (\texttt{eclipse-sumo}, \texttt{libsumo}),
        Stable-Baselines3, Gymnasium, PyTorch. Exact versions are in \texttt{requirements.txt}.
  \item \textbf{Hardware:} Kaggle notebooks with 4 CPU cores. A GPU is optional; a benchmark on a
        Tesla T4 gave 62.9\,s against 64.0\,s on the CPU for 6{,}000 training steps, because the
        simulator, not the network, takes most of the time.
  \item \textbf{Seeds:} training seeds 1 to 5; route seeds 0, 1 and 2 for each test day.
  \item \textbf{Commands} (from \texttt{src/}):
\begin{verbatim}
python -m data.download
python -m data.preprocess
python train.py --delta-t 2 --warmstart 20000 --seed 1 --name dqn_camera_seed1
python evaluate.py --models ../results/dqn_camera_seed1.zip
python plots.py
python report_tables.py
\end{verbatim}
  \item \textbf{Notebook:} \texttt{notebooks/pune\_signal\_rl.ipynb} runs the full pipeline on
        Kaggle or Google Colab.
\end{itemize}

\clearpage
\nocite{*}
\begin{thebibliography}{11}
\bibitem{mnih2015} V. Mnih et al., ``Human-level control through deep reinforcement learning,''
  \emph{Nature}, vol. 518, pp. 529--533, 2015.
\bibitem{sutton2018} R. S. Sutton and A. G. Barto, \emph{Reinforcement Learning: An Introduction},
  2nd ed., MIT Press, 2018.
\bibitem{pune2026} A. Matange and J. Abraham, ``Heterogeneous Traffic Count Dataset, Pune (Jan 2023),''
  Mendeley Data, V1, 2026. doi:10.17632/xnf2k6n288.1.
\bibitem{irc106} Indian Roads Congress, \emph{IRC:106-1990 Guidelines for Capacity of Urban Roads in
  Plain Areas}, New Delhi, 1990.
\bibitem{sumo} P. A. Lopez et al., ``Microscopic traffic simulation using SUMO,''
  in \emph{Proc. IEEE ITSC}, 2018, pp. 2575--2582.
\bibitem{sb3} A. Raffin et al., ``Stable-Baselines3: Reliable reinforcement learning implementations,''
  \emph{JMLR}, vol. 22, no. 268, pp. 1--8, 2021.
\bibitem{wei2019} H. Wei, G. Zheng, V. Gayah and Z. Li, ``A survey on traffic signal control methods,''
  arXiv:1904.08117, 2019.
\bibitem{genders2016} W. Genders and S. Razavi, ``Using a deep reinforcement learning agent for
  traffic signal control,'' arXiv:1611.01142, 2016.
\bibitem{wei2018} H. Wei, G. Zheng, H. Yao and Z. Li, ``IntelliLight: A reinforcement learning approach
  for intelligent traffic light control,'' in \emph{Proc. ACM SIGKDD}, 2018, pp. 2496--2505.
\bibitem{dabney2018} W. Dabney, M. Rowland, M. G. Bellemare and R. Munos, ``Distributional reinforcement
  learning with quantile regression,'' in \emph{Proc. AAAI}, 2018, pp. 2892--2901.
\end{thebibliography}
\end{document}
